%% ma: L11-B16-0038
%% lop: 11
%% chuong: 5
%% bai: 16
%% dang: 11.16.4
%% loai: TN4
%% muc: VD
%% dap_an: "A"
%% nguon: "tai-lieu-hoc-tap-mon-toan-lop-11-theo-chuong-trinh-giao-duc-pho-thong-2018.pdf (Lớp Toán Thầy Duy, Phạm Lê Duy, THPT Chu Văn An), Chương 3 Bài 2, Phần 1, Câu 26 (tr.366)"
%% kien_thuc: "Lớp 11 Bài 16 (giới hạn của hàm số)"
%% trang_thai: cho-duyet
%% ly_do: "Dạng mới đề xuất (11.16.x), chờ giáo viên duyệt dạng; tài liệu gốc không có lời giải mẫu, đáp án do kho tự giải và kiểm bằng sympy."
\begin{ex}
Biết rằng $\lim\limits_{x\to+\infty}\left(\dfrac{x^2+1}{x-2}+ax-b\right)=-5$. Tính tổng $a+b$.
\choice{\True $6$}{$7$}{$8$}{$5$}
\loigiai{\sloppy $\dfrac{x^2+1}{x-2}=x+2+\dfrac5{x-2}$ nên biểu thức $=(1+a)x+(2-b)+\dfrac5{x-2}$. Suy ra $1+a=0$ và $2-b=-5$: $a=-1$, $b=7$. Vậy $a+b=6$.}
\end{ex}
